%% GENERATED by python/paper/build.py from results/ -- do not edit by hand.
%% Rebuild:  .venv\Scripts\python.exe python\paper\build.py --only US_PensChars
%% Source:   results/shocks/US_shocksCommonX.csv
\begin{table}[!htb]
\centering
\begin{threeparttable}
\caption{The effect of pension design ($\theta$) in US -- 2020}
\label{table:US:pensChars}
\renewcommand{\arraystretch}{1.25}
\begin{tabularx}{.9\textwidth}{p{3cm}YYY}
\toprule
\textbf{Scenario} & \textbf{Tax rate} & \textbf{Savings rate} & \textbf{Avg. workweek} \\
\midrule 
$\theta = 0.74$ & 14.43\% & 15.37\% & 39.39 \\[1.25ex]
\multicolumn{4}{l}{\textit{Full effect:}} \\\hline
$\theta = 0$ & 18.52\% & $-0.93$ p.p. & 37.52 \\
$\theta = 1$ & 12.28\% & $+0.65$ p.p. & 40.13 \\[1.25ex]
\multicolumn{4}{l}{\textit{Economic equilibrium effect:}} \\\hline
$\theta = 0$ & 14.43\% & $+0.41$ p.p. & 38.54 \\
$\theta = 1$ & 14.43\% & $-0.14$ p.p. & 39.67 \\[1.25ex]
\bottomrule
\end{tabularx}
\begin{tablenotes}
\footnotesize
\item $\rho = 1.0$. The economic-equilibrium rows hold $\tau$ at the baseline path, so they isolate the response of savings and hours to $\theta$ alone; the full rows re-optimise $\tau$ politically. The savings rate is savings relative to GDP; the baseline row reports its level and every other row the change against that baseline in percentage points. Common-$X$ calibration: one leisure parameter $X$ shared across income groups, with the hours unit pinned by targeting the observed average workweek, so relative hours are a prediction rather than a calibration target.
\end{tablenotes}
\end{threeparttable}
\end{table}
